% !TEX program = xelatex
\documentclass[12pt]{article}
\usepackage[margin=28mm]{geometry}
\usepackage{hyperref}
\usepackage{amsmath, amsthm, amssymb, amsfonts}
\usepackage{tikz-cd}
\usepackage{fontspec}
\usepackage{unicode-math}
\setmainfont{Latin Modern Roman}
\setmonofont{Latin Modern Mono}
\usepackage{listings}
\lstset{basicstyle=\ttfamily\footnotesize, columns=fullflexible, breaklines=true, showstringspaces=false}

\title{lean\_basic\_task — Lean→TeX Export (English)}
\author{TeX generated by ChatGPT-5 Pro\\Lean source generated by CODEX (OpenAI)}
\date{2025-09-10}

\begin{document}
\maketitle

\begin{abstract}
This document was automatically generated from a Lean source file. It includes provenance, sample diagrams, a quick stats table, a declarations table (first-line signatures), and the full original Lean listing.
\end{abstract}

\section*{Provenance}
This \LaTeX{} file was generated from the Lean source file \texttt{lean\_basic\_task.lean}.\\
\textbf{TeX generation}: ChatGPT-5 Pro.\\
\textbf{Lean source generation}: CODEX (OpenAI).

\section*{At-a-glance}
Total lines: 25\\
Defs: 0\\
Lemmas: 0\\
Theorems: 2\\
Structures: 0\\
Inductives: 0

\section*{Sample diagrams}

\[
\begin{tikzcd}
A \arrow[r, "f"] \arrow[d, "g"'] & B \arrow[d, "h"] \\
C \arrow[r, "k"'] & D
\end{tikzcd}
\]



\[
\begin{tikzcd}
T \arrow[r, "h"] \arrow[dr, swap, "\pi_1\circ h"] \arrow[drr, "\pi_2\circ h"] & X\times Y \arrow[d, swap, "\pi_1"] \arrow[dr, "\pi_2"] \\
& X & Y
\end{tikzcd}
\]


\section*{Declarations (first-line signatures)} 
\begin{center}
\begin{tabular}{lll}
\textbf{kind} & \textbf{name} & \textbf{signature (first line)} \\ \hline
theorem & \texttt{zero\_scalar\_mul} & \texttt{{R V : Type*} [Semiring R] [AddCommMonoid V] [Module R V] (v : V) : (0 : R) • v = 0 := by} \\
theorem & \texttt{scalar\_mul\_zero} & \texttt{{R V : Type*} [Semiring R] [AddCommMonoid V] [Module R V] (a : R) : a • (0 : V) = 0 := by} \\
\end{tabular}
\end{center}

\section*{Lean source (verbatim)}
\begin{lstlisting}
import Mathlib.Algebra.Module.Basic

-- ゼロスカラーをかけるとゼロベクトルになる
theorem zero_scalar_mul {R V : Type*} 
  [Semiring R] [AddCommMonoid V] [Module R V]
  (v : V) : 
  (0 : R) • v = 0 := by
  simpa using (zero_smul R v)
  
-- ゼロベクトルに何をかけてもゼロベクトル
theorem scalar_mul_zero {R V : Type*} 
  [Semiring R] [AddCommMonoid V] [Module R V]
  (a : R) : 
  a • (0 : V) = 0 := by
  simpa using (smul_zero a)

-- Bourbaki スタイルの軽い example（API の最小行使）
section Examples
  variable {R V : Type*}
  variable [Semiring R] [AddCommMonoid V] [Module R V]
  example (v : V) : (0 : R) • v = 0 := by
    simpa using (zero_smul R v)
  example (a : R) : a • (0 : V) = 0 := by
    simpa using (smul_zero a)
end Examples

\end{lstlisting}

\end{document}
